\paragraph{C.3p (Known with derivation): local etale Lefschetz, Picard torsion, and completion.}
Let $(R,\mathfrak m)$ be an excellent normal strictly
henselian local ring, let $0\ne f\in\mathfrak m$,
and write
\[
 X=\operatorname{Spec}R\setminus\{\mathfrak m\},
 \qquad Y=X\cap V(f),\qquad N=\dim R.
\]
Assume $X$ regular. Restriction of finite etale covers
is fully faithful for $N\ge3$ and is an equivalence
for $N\ge4$:
\begin{equation}\label{eq:C3-etale-Lefschetz}
 \operatorname{FEt}(X)\xrightarrow{\ \sim\ }
 \operatorname{FEt}(Y)\qquad(N\ge4).
\end{equation}
If also $Y$ is regular, then, for every positive
integer $n$ invertible in $R$,
\begin{equation}\label{eq:C3-Picard-torsion-Lefschetz}
 \operatorname{Pic}(X)[n]\xrightarrow{\ \sim\ }
 \operatorname{Pic}(Y)[n]\qquad(N\ge4).
\end{equation}
This assertion includes mixed characteristic and
has no coefficient-field or algebraizability
hypothesis.

Separately, in the coefficient-field setting of
C.3n, completion induces a canonical isomorphism
of the represented local Picard functors
\begin{equation}\label{eq:C3-Picard-completion}
 \mathbf{Pic}^{\mathrm{loc}}_{R/k}
       \xrightarrow{\ \sim\ }
 \mathbf{Pic}^{\mathrm{loc}}_{\widehat R/k}
                     \qquad(N\ge3).
\end{equation}
The functors in this last assertion use
henselized tensor products as defined in C.3n.
It does not identify their arbitrary parameter
families with line bundles on the fixed product
$X\times T$.

\begin{proof}[Proof of the finite-etale and torsion assertions]
Put $A=\widehat R$ and use hats for the base changes
of $X,Y$. These are their punctured completions:
the closed fibre of $\operatorname{Spec}A\to
\operatorname{Spec}R$ has just the one closed
point. Excellence makes $A$ a normal complete
local domain of dimension $N$. It also makes
$R\to A$ regular. Thus $\widehat X$ is regular:
each point above the puncture is obtained by
localizing a regular morphism over a regular
local ring. The completed ring has a dualizing
complex.

Elkik's exact input is
\citet[Section III, Theorem 5, pp.~577--579, and
its corollary on p.~579]{D_Elkik1973}.
For a noetherian henselian pair $(B,I)$, completion
induces an equivalence of the categories of
finite etale covers of
$\operatorname{Spec}B\setminus V(I)$ and its
completed counterpart. Apply this both to
$(R,\mathfrak m)$ and to $(R/(f),\mathfrak m/(f))$.
The resulting vertical functors in
\[
\begin{array}{ccc}
 \operatorname{FEt}(X)&\longrightarrow&\operatorname{FEt}(Y)\\
 \big\downarrow\rlap{$\scriptstyle\sim$}&&
 \big\downarrow\rlap{$\scriptstyle\sim$}\\
 \operatorname{FEt}(\widehat X)&\longrightarrow&
 \operatorname{FEt}(\widehat Y)
\end{array}
\]
are equivalences. The hypotheses here concern
noetherian henselian pairs; no principal-ideal
restriction from Elkik's later Theorem 7 is
being added or silently removed.

It remains to check the complete-ring restriction.
The ring $A$ is $f$-adically complete. Indeed
an $f$-adically Cauchy sequence is
$\mathfrak m_A$-adically Cauchy, and each
ideal $(f^r)$ is closed for the latter topology
by the Krull intersection theorem applied to
$A/(f^r)$. Its limit therefore has the required
residues modulo every $f^r$.
The ring $A_f$ is regular and hence satisfies
$(S_2)$. For $N\ge3$, all the hypotheses of
\citet[Lemma 58.19.5, Tag 0EK5]{D_StacksEtaleLefschetz}
hold, giving full faithfulness.

For essential surjectivity with $N\ge4$, apply
\citet[Lemma 58.24.1, Tag 0EK9]{D_StacksEtaleLefschetz}.
Its additional condition is purity at every
maximal ideal of $A_f$. Here is the dimension
check for that condition. The punctured spectrum
of a noetherian local ring is Jacobson
\citep[Lemma 28.6.4, Tag 02IM]{D_StacksEtaleLefschetz}.
Hence a point closed in the open $\operatorname{Spec}A_f$
is closed in the puncture. If its prime in $A$
is $\mathfrak p$, this says
$\dim A/\mathfrak p=1$. The complete local domain
$A$ is catenary and equidimensional, so
\[
                \dim A_{\mathfrak p}=N-1\ge3.
\]
This is a regular local ring of dimension at
least two. Purity therefore holds by
\citet[Lemma 58.21.3(1), Tag 0BMA]{D_StacksEtaleLefschetz}.
This is the categorical purity clause (1);
no assertion about finite algebras with extra
components supported at the closed point is
needed. Lemma 58.24.1 now proves essential
surjectivity. Together with full faithfulness
and the Elkik comparisons it proves
\eqref{eq:C3-etale-Lefschetz}.

For \eqref{eq:C3-Picard-torsion-Lefschetz}, use
the Kummer exact sequence with $n$ invertible.
The global units on $X$ are $R^\times$ by
normal Hartogs extension, and every such unit
has an $n$th root in $R$ by strict henselianity.
For $Y$, use the finite normalization from
C.3k(i): the normal strictly henselian factors
$C_a$ have dimension $N-1\ge3$, their punctures
are the disjoint components of $Y$, and
\[
                 \Gamma(Y,\mathcal O_Y)=\prod_a C_a.
\]
Their units likewise have $n$th roots.
Consequently Kummer gives natural isomorphisms
\[
 H^1_{\mathrm{et}}(X,\mu_n)=\operatorname{Pic}(X)[n],
 \qquad
 H^1_{\mathrm{et}}(Y,\mu_n)=\operatorname{Pic}(Y)[n].
\]
The finite-etale equivalence preserves torsors
under the finite etale group $\mu_n$ and their
group actions. It thus gives the required
isomorphism. In dimension three the same
argument with full faithfulness gives the
injectivity already proved independently in
C.3k(i).
\end{proof}

\begin{proof}[Separate proof of the local Picard functor comparison]
Let $B$ be a finite-type $k$-algebra and form the
two henselian pairs
\[
 S_B=(R\otimes_k B)^h_{\mathfrak m},\qquad
 T_B=(\widehat R\otimes_k B)^h_{\mathfrak m\widehat R}.
\]
Both rings are noetherian: their unhenselized
rings are finite type over noetherian rings,
and henselization of a noetherian pair is
noetherian. The natural map $S_B\to T_B$
induces the canonical identification of their
completions along the displayed ideals,
\begin{equation}\label{eq:C3-common-tensor-completion}
 \widehat S_B=\widehat T_B
       =\varprojlim_r (R/\mathfrak m^r)\otimes_k B.
\end{equation}
In particular the completed complements of
these ideals are the same scheme.

The exact source now is
\citet[Section III, Theorem 3, p.~569, and its
unnumbered Picard corollary on p.~573]{D_Elkik1973}:
for every noetherian henselian pair, the Picard
group of the ideal complement maps bijectively
to that of its completion. Apply it to $S_B$
and $T_B$. Both Picard groups map bijectively
to the Picard group of the common completed
complement in \eqref{eq:C3-common-tensor-completion}.
The natural map between them is therefore
a bijection. These maps are functorial in $B$,
so the two local Picard presheaves, and then
their etale sheafifications, agree on finite-type
$k$-algebras.

Boutot's representability theorem applies to
both $R$ and $\widehat R$ in the hypotheses
of C.3n. The resulting group schemes are
locally of finite type over $k$. The canonical
morphism between them is an isomorphism
because it is a natural bijection on every
finite-type $k$-algebra: the affine finite-type
open subschemes of these group schemes test
the inverse by Yoneda, and the inverses glue
by their uniqueness. Equivalently one may
extend from finite-type parameter algebras
by their filtered-colimit property.
This proves \eqref{eq:C3-Picard-completion}.

Every comparison in this proof occurs after
henselizing the tensor pair. None removes
that henselization from an arbitrary family
on a proper parameter scheme. Therefore
\eqref{eq:C3-Picard-completion}, even combined
with \eqref{eq:C3-Picard-torsion-Lefschetz},
does not by itself prove surjectivity for
ordinary algebraic equivalence on fixed
products. This precise remaining distinction
is compatible with the characteristic-zero
counterexample C.3m.
\end{proof}
